\documentclass{article}
\usepackage{amsmath}
\begin{document}
\title{Earliest robust commitment for a source-conditioned interruption}
\maketitle

\paragraph{Kernel.}
One authorized actor must choose one irreversible release time.
Before releasing, its legitimate observation stream is the
projection of a history under the policy that has not yet released.
No other probing or resource action is selectable by this actor.
Release can change the future history; the model must already
supply the complete quantified acceptance predicate
\(A(h,t)\) for releasing at \(t\) in source history \(h\),
under the specified residual response law and all admitted
post-release outcomes. This formulation neither supplies those
controllers nor assumes the physical effects are readable. The
set of candidate histories is fixed independently of success and
includes every admitted pre-release response. Time ticks are
available to the actor; arbitrary time uncertainty may instead be
included in the history if its actual observation is represented.
Policies are deterministic and must eventually release in each
candidate history.

\paragraph{Exact criterion and earliest rule.}
Let \([h]_t\) be the histories agreeing with \(h\) on the
actor's entire observed prefix through \(t\), including negative
observations where justified. A robust causal release policy
exists exactly when
\[
 \forall h\;\exists t:\; \forall h'\in[h]_t,\ A(h',t).
\]
Necessity fixes \(t\) to the policy's action in \(h\).
Every history in the actor's same prefix class must take the same
action and satisfy the release predicate. For sufficiency, let
the actor release at the first cut where all compatible
histories permit release. The time of the earliest such cut is
the same for histories with equal observed prefix at that cut,
since their earlier prefixes are also equal.
Thus this is a stopping rule and each chosen action is accepted.
It releases no later on any history than any other winning rule.
The Lean statement uses classical selection to express the least
tick and proves its causality and optimality; for a finite listed
support with decidable acceptance, the implementation scans the
attained prefix classes using only the current local transcript.
No computability of arbitrary physical acceptance sets is inferred.

A criterion that only checks each world's feasible cuts can
accept a family which no single actor can realize. Conversely a
test that intersects all histories once, at time zero, can
reject a family that a future observation separates in time.
If each acceptance set is a closed time interval or empty, each
nonempty attained class reduces to the largest of its lower
endpoints and the smallest of its upper endpoints; empty
case-specific intervals must remain explicit. The relevant class
cannot in general be computed from the latest value alone.

\paragraph{Two source interpretations.}
The staged ARIAC replay (\#778) includes old floor work, a
public camera, and a release that both introduces a new duty
and shuts off the public feed. The capture evidence is retained
before release under an explicit permission; the pose must be
valid after release in the actual parking/authority context.
A shift that breaks this validity can make \(A\) false even
when the observation of the old context was truthful. The
physical capture seam and the conditional source model are
separate from the proof in this note.

The staged SubT replay (\#779) keeps the old receipt alive
until a near relay's route starts. In the scoped known-history
case, native dispatch and registered route determine exactly
\(p<k\le b\) for effective registration \(k\).
Positive bounded activation latency lifts this to
\(p<s+L,\ s+U\le b\) for command \(s\).
The independent authority to command is a premise.
If the old receiving callback reaches the caller at \(r\) and
the caller may wait until then, feasibility of this restricted
receipt-first policy reduces to \(r+U\le b\).
In the two named histories a local receipt separates the
critical cuts, while a one-tick-later remote report does
not. An available far relay changes the actual predicate, not
the observed prefix of the near-relay policy.

\paragraph{Preserved distinctions.}
The common rule describes the timing of an irreversible act.
It does not identify the assembly pose with a transmitted
receipt, derive release authority from having data, or claim
new source observability. A substitution may preserve
physical acceptability at all cuts and still remove an early
prefix distinction the policy requires. The proof gives no
conclusion if an alternate probe changes the pre-commit
history or if several local actors may each commit. This is a
verified finite-observation/stopping specialization.
Direct causal policy enumeration agrees; no comparative
performance, independent transfer, or architecture completion
is asserted.
\end{document}
